\begin{textsample}{POS Dim 1 – vem\_ed\_08 – Score 6.00 – vem\_ed\_08/t029.txt}  \label{ex:f1_pos_068}
QUEBRA-GELO A equipe dos Projetos Colaborativos Internacionais (PCIs / Cesu) vem trabalhando para o aprimoramento \textbf{contínuo} dos PCIs nas Fatecs.
E também contribui para a \textbf{expansão} dos Intercâmbios Virtuais para o Ensino Médio e Técnico.
Após um contato da equipe dos PCIs com a Universidad de Monterrey (UDEM), no México, surgiu a ideia de levar a expertise dos PCIs desenvolvidos nas Fatecs também para as Etecs.
Com a orientação da equipe dos PCIs / Cesu, a Cetec (Unidade de Ensino Médio e Técnico) elaborou com a ArInter (Assessoria de Relações Internacionais do CPS) as linhas \textbf{gerais} do Programa de Aprendizagem Colaborativa Internacional.
Em abril de 2021, a equipe dos PCIs conduziu o projeto piloto entre Etec Jaraguá e PREPA UDEM (escolas preparatórias para a universidade).
Uma notícia do site Cesu relata essa iniciativa: https://cesu.cps.sp.gov.br/arrie sguemos-el-portunhol-expandindo-os-intercambios-virtuais-para-o-ensino-medio/.
Olhando os PCIs nas Fatecs, os \textbf{números} do primeiro semestre de 2021 \textbf{são} animadores: 725 estudantes participaram de Intercâmbios Virtuais.
Ao \textbf{todo}, 30 professores de 20 Fatecs participaram de PCIs com 14 IES internacionais.
Pesquisa realizada em julho com estudantes e professores participantes dos PCIs demonstra \textbf{que} a satisfação com essa experiência \textbf{é} elevada.
A síntese dos principais resultados colhidos junto aos alunos está na página 4.
Na próxima edição, iremos \textbf{publicar} o resumo das percepções de docentes.
As páginas 2 e 3 trazem uma entrevista com a vice-diretora-superintendente do CPS, Emilena Lorenzon Bianco.
Na seção “Boas Práticas”, o relato de Maira Rezende \textbf{destaca} \textbf{que} a “comunicação entre os alunos e a organização do tempo” estão entre os principais desafios ao realizar um PCI.
Vale a pena encará-los pois “a experiência [... ] \textbf{é} incrível”.
Boa leitura!

% matched lemmas: contínuo, destacar, expansão, geral, número, publicar, que, ser, todo
\end{textsample}
